%=====================================================================
\chapter{Regression and Modelling for Inference}\label{ch:regression}
%=====================================================================
\locator{4}{Part~4 --- Measurement, Data and Analysis}

\begin{outcomes}
By the end of this chapter you will be able to:
\begin{tight}
  \item \textbf{Interpret} a regression coefficient correctly, including the
    conditioning clause everyone omits.
  \item \textbf{Choose} among linear, logistic, Poisson and negative binomial
    models by the shape of the outcome.
  \item \textbf{Specify} a mixed-effects model when observations are nested.
  \item \textbf{Distinguish} modelling for explanation from modelling for
    prediction, and stop using the tools of one for the other.
  \item \textbf{Report} a model completely enough to be reproduced.
\end{tight}
\end{outcomes}

\begin{prereq}
\chref{inference} for assumptions and intervals; \chref{questions} for
variable roles --- a regression coefficient means something quite different
depending on whether the covariate is a confounder or a mediator.
\end{prereq}

\begin{keyterms}
linear model $\bullet$ coefficient $\bullet$ residual $\bullet$ logistic
regression $\bullet$ odds ratio $\bullet$ Poisson and negative binomial
$\bullet$ overdispersion $\bullet$ mixed-effects model $\bullet$ random
intercept and slope $\bullet$ collinearity $\bullet$ VIF $\bullet$ stepwise
selection $\bullet$ explanation versus prediction
\end{keyterms}

\section{What a Coefficient Means}

\begin{definitionbox}[The conditioning clause]
In $y = \beta_0 + \beta_1 x_1 + \beta_2 x_2 + \varepsilon$, the coefficient
$\beta_1$ is the expected change in $y$ for a one-unit increase in $x_1$
\textbf{holding $x_2$ fixed}.

\medskip
That clause is not decoration. It means $\beta_1$ answers a different question
depending on what else is in the model, and adding or removing a covariate
changes its value and sometimes its sign. There is no such thing as ``the
effect of $x_1$'' independent of the model it sits in.
\end{definitionbox}

\begin{alertbox}[Which covariates to include is a causal question, not a statistical one]
The habit of including every available variable ``to control for them'' is
actively harmful, and which variables belong depends on their causal role
(\chref{questions}).

\begin{tight}
  \item \textbf{Confounders} --- causes of both treatment and outcome ---
    \emph{must} be included, or the estimate is biased.
  \item \textbf{Mediators} --- on the causal path from treatment to outcome ---
    must \emph{not} be included if you want the total effect. Controlling for a
    mediator removes exactly the effect you were trying to measure.
  \item \textbf{Colliders} --- common effects of treatment and outcome ---
    must not be included: conditioning on one \emph{creates} an association
    where none existed.
\end{tight}

\medskip
Draw the assumed causal structure before choosing the model. ``I controlled for
everything I had'' is not a defence; it is a description of a model whose
coefficients cannot be interpreted.
\end{alertbox}

\section{Choosing a Model by the Shape of the Outcome}

\begin{center}
\small
\begin{longtable}{@{}L{3.0cm}L{3.6cm}L{7.0cm}@{}}
\toprule
\textbf{Outcome} & \textbf{Model} & \textbf{Coefficient means} \\
\midrule
\endfirsthead
\toprule
\textbf{Outcome} & \textbf{Model} & \textbf{Coefficient means} \\
\midrule
\endhead
\bottomrule
\endfoot
Continuous, roughly symmetric
  & Linear regression
  & Change in $y$ per unit of $x$, in $y$'s units \\
Binary
  & Logistic regression
  & Change in \emph{log odds}. Exponentiate for an odds ratio --- which is
    \emph{not} a risk ratio, and is routinely misread as one \\
Count, variance $\approx$ mean
  & Poisson regression
  & Multiplicative change in the expected count, after exponentiating \\
Count, variance $>$ mean
  & Negative binomial
  & As Poisson. \textbf{This is the usual case in software data} --- defects,
    commits and issues are almost always overdispersed \\
Count with excess zeros
  & Zero-inflated or hurdle model
  & Two processes: whether any occur, and how many given some do \\
Time to event
  & Cox proportional hazards
  & Hazard ratio; handles censoring, which matters for time-to-fix data \\
Ordinal
  & Ordinal (proportional odds) regression
  & Respects ordering without assuming equal spacing
    (\chref{measurement}) \\
\end{longtable}
\end{center}

\begin{pitfall}[Fitting Poisson to overdispersed counts]
Defect counts, commit counts and issue counts nearly always have variance
substantially greater than their mean, because the units are heterogeneous ---
some modules are simply more defect-prone than others.

\medskip
Fitting a Poisson model to such data does not produce wrong coefficients so
much as \emph{far too narrow} standard errors. The result is confidence
intervals that are too tight and p-values that are too small, and the entire
inference is wrong in the direction of over-confidence.

\medskip
Check: compute the residual deviance divided by its degrees of freedom. If it
is much above 1, use a negative binomial model. In R this is one function call
--- \texttt{MASS::glm.nb} in place of \texttt{glm(family = poisson)}.
\end{pitfall}

\section{Nested Data and Mixed-Effects Models}

\chref{experiments} warned that analysing 200 reviews from 20 developers as
$N = 200$ inflates significance. The mixed-effects model is the principled fix.

\begin{center}
\small
\begin{tabular}{@{}L{3.4cm}L{10.9cm}@{}}
\toprule
\textbf{Fixed effect} & The population-level effect you want to estimate ---
                        the treatment \\
\textbf{Random intercept} & Each developer has their own baseline. Removes the
                        assumption that all developers are alike \\
\textbf{Random slope}  & The treatment effect itself differs by developer. Fit
                        this when it is plausible, which it usually is \\
\bottomrule
\end{tabular}
\end{center}

\begin{lstlisting}[language=Rlang, caption={A mixed-effects model for the code review
experiment, with reviewer and the reviewed artefact as crossed random effects.}]
library(lme4); library(lmerTest)

# Reviewers and artefacts are crossed, not nested: every reviewer sees
# several artefacts and every artefact is seen by several reviewers.
m <- lmer(defects_found ~ technique + experience +
            (1 + technique | reviewer) + (1 | artefact),
          data = reviews)

summary(m)
confint(m, method = "profile")     # intervals, not just p-values

# Does the random slope earn its place?
m0 <- lmer(defects_found ~ technique + experience +
             (1 | reviewer) + (1 | artefact), data = reviews)
anova(m0, m)
\end{lstlisting}

\begin{conceptbox}[When you must use a mixed model]
Whenever observations share something that makes them more alike than
observations from different groups:

\begin{tight}
  \item Several measurements from the same participant.
  \item Several files from the same repository.
  \item Several students from the same class or supervisor.
  \item Repeated measures over time on the same unit.
\end{tight}

\medskip
Ignoring this structure is one of the most consequential errors available, and
it is invisible in the output: the model runs, the coefficients look fine, and
the standard errors are wrong. Decide the random-effects structure from the
\emph{design}, before fitting --- it is a description of how the data was
generated, not a modelling choice to be selected by fit
statistics~\cite{bates2015fitting}.
\end{conceptbox}

\section{Explanation Is Not Prediction}

\begin{paperbox}[The distinction that resolves most modelling arguments]
Galit Shmueli, \emph{To Explain or to Predict?}, Statistical Science 25(3),
2010~\cite{shmueli2010explain}.\oa{}

\medskip
\textbf{Why it matters.} Explanatory and predictive modelling have different
goals, and therefore different criteria, different variable selection, and
different definitions of a good model. Confusing them is behind a great deal of
muddled work --- especially in computing, where scholars trained on machine
learning apply its habits to inferential questions.

\medskip
\textbf{What to extract.} That a model can be excellent for prediction and
useless for explanation, and the reverse. Decide which you are doing, in
writing, before you fit anything.
\end{paperbox}

\begin{center}
\small
\begin{tabular}{@{}L{3.0cm}L{5.4cm}L{5.9cm}@{}}
\toprule
 & \textbf{Explanation} & \textbf{Prediction} \\
\midrule
Goal        & Estimate an effect and its uncertainty
            & Minimise error on new data \\
Variables   & Chosen by causal role
            & Chosen by predictive contribution; causality irrelevant \\
Judged by   & Unbiasedness; correct standard errors
            & Held-out performance (\chref{mleval}) \\
Collinearity& A serious problem: it destroys the interpretability of individual
              coefficients
            & Largely harmless: predictions stay accurate \\
Overfitting & Less central, though still real
            & The central concern \\
Report      & Coefficients with intervals
            & Performance on a test set, with the validation scheme \\
\bottomrule
\end{tabular}
\end{center}

\begin{alertbox}[Stepwise selection has no place in explanatory modelling]
Automated forward, backward or stepwise selection on p-values produces biased
coefficients, standard errors that are far too small, p-values that are
meaningless, and $R^2$ that is inflated. It also gives unstable results ---
resample the data and a different model is selected.

\medskip
The reason is simple: the selection process is itself an analysis of the data,
and the reported inference does not account for it. This is the garden of
forking paths (\chref{inference}) automated.

\medskip
For explanation, specify the model from the causal structure and report it. For
prediction, use regularisation --- lasso, ridge, elastic net --- with
performance judged on held-out data, and do not interpret the resulting
coefficients as effects.
\end{alertbox}

\section{Collinearity}

When predictors are strongly correlated, the model cannot separate their
contributions. Coefficients become unstable and their standard errors inflate,
even though overall fit is unaffected.

\begin{center}
\small
\begin{tabular}{@{}L{3.6cm}L{10.7cm}@{}}
\toprule
Diagnose   & Variance inflation factor. Above 5 warrants attention; above 10
             usually indicates a real problem \\
Accept it  & If the collinear variables are controls and you do not interpret
             their coefficients, collinearity among them is not a problem \\
Combine    & If two predictors measure the same construct, combine them into
             one scale (\chref{measurement}) rather than fighting the model \\
Drop       & Only with a substantive reason. Dropping the one with the larger
             p-value is a data-driven decision that invalidates the inference \\
\bottomrule
\end{tabular}
\end{center}

Note the classic case in software engineering: cyclomatic complexity and lines
of code correlate strongly. A model containing both cannot tell you which
predicts defects, and a paper claiming complexity predicts defects while
controlling for size is making a much weaker claim than it sounds.

\section{Reporting a Model}

\begin{center}
\small
\begin{tabular}{@{}L{4.0cm}L{10.3cm}@{}}
\toprule
The full specification & Every term, including random effects, in a form that
                         could be retyped \\
Why these covariates   & The causal reasoning, not ``available in the dataset'' \\
Coefficients           & With confidence intervals, in interpretable units.
                         Exponentiate log-odds and count-model coefficients \\
Model checks           & Residual plots, overdispersion, VIF, influential points \\
Sample size            & At every level: how many observations, how many
                         developers, how many repositories \\
Missing data           & How much, and how handled. Listwise deletion is a
                         choice with consequences, not a neutral default \\
Software               & Package and version; the model syntax verbatim
                         (\chref{reproducibility}) \\
\bottomrule
\end{tabular}
\end{center}

\begin{pitfall}[Three interpretation errors that reach print]
\textbf{Reading an odds ratio as a risk ratio.} An odds ratio of 2.0 does not
mean the outcome is twice as likely. When the outcome is common the two diverge
sharply, and the odds ratio always looks more dramatic.

\medskip
\textbf{Interpreting a coefficient as an effect in observational data.}
Regression adjusts for what you measured; it does not create the counterfactual
(\chref{quasiexp}).

\medskip
\textbf{Comparing standardised coefficients across models or samples.}
Standardisation divides by that sample's standard deviation, so a
``standardised'' coefficient is not comparable across populations with
different variances.
\end{pitfall}

\begin{traditions}{regression modelling}
\tradEMP{The workhorse. Model specification follows from the causal structure,
and the report is coefficients with intervals plus diagnostics.}
\tradDSR{Used to model the relationship between artefact configuration and
outcome, and to identify which design parameters matter.}
\tradQUAL{Not used. The nearest analogue --- explaining variation by tracing
mechanisms across cases --- is \chref{casestudy}'s explanation building.}
\tradFORM{Regression models have properties that are proved: identifiability,
consistency, asymptotic normality. Knowing which of these fail under your data
is what distinguishes fitting from modelling.}
\end{traditions}

\begin{lab}[Fit, then break, a model]
Take a dataset with a nested structure --- multiple observations per unit.

\begin{tightnum}
  \item Draw the assumed causal structure, and mark each variable as confounder,
    mediator or collider.
  \item Fit an ordinary regression ignoring the nesting. Record the standard
    error on your main coefficient.
  \item Fit the mixed-effects model with the correct random structure. Record
    the standard error again.
  \item Compare. The ratio is how much the naive model overstated your
    precision.
  \item Now add a mediator to the model and watch the main coefficient move.
    Explain in one sentence why it moved.
\end{tightnum}
\end{lab}

\begin{checkpoint}
\begin{tightnum}
  \item Why does the meaning of $\beta_1$ depend on what else is in the model?
  \item Your outcome is a defect count with variance four times its mean. Which
    model, and what goes wrong if you use Poisson?
  \item When must you fit a random slope rather than only a random intercept?
  \item Why is stepwise selection unacceptable for explanation but tolerable
    for prediction?
\end{tightnum}
\end{checkpoint}

\begin{chaptersummary}
\begin{tight}
  \item A coefficient means change in the outcome \emph{holding the other terms
    fixed}. Which covariates to include is a causal question: include
    confounders, exclude mediators and colliders.
  \item Choose the model from the outcome's shape. Software counts are
    overdispersed; use negative binomial, not Poisson.
  \item Nested observations require a mixed-effects model, and the random
    structure follows from the design, not from fit statistics.
  \item Explanation and prediction have different goals, criteria and variable
    selection. Decide which you are doing before fitting.
  \item Stepwise selection invalidates inference. Specify from causal structure,
    or regularise and do not interpret coefficients.
  \item Report the full specification, intervals in interpretable units,
    diagnostics, sample sizes at every level, missing data handling, and
    software versions.
\end{tight}
\end{chaptersummary}

\begin{reviewq}
  \item Controlling for a mediator removes the effect you wanted to measure,
    yet ``controlling for more variables'' is widely treated as more rigorous.
    Explain the misconception and give a software engineering example where it
    would reverse a conclusion.
  \item Shmueli argues explanation and prediction require different models.
    Computing research increasingly uses machine learning methods for
    inferential questions. What specifically goes wrong, and how would you
    detect it as a reviewer?
  \item Complexity metrics correlate strongly with size. Design a study that
    could establish whether complexity predicts defects independently of size,
    and say what data it would require.
  \item Mixed-effects models are standard in psychology and rare in software
    engineering, despite our data being at least as nested. Account for the
    gap, and say what would close it.
\end{reviewq}
